% ============================================================================
% Fichier  : partie4/P04_C14_cas_integrateurs.tex
% Titre    : Études de Cas Intégratrices
% Auteur   : MINKA MI NGUIDJOÏ Thierry Emmanuel
% Version  : 1.0 -- novembre 2026
% Corrigés : corriges/Corr_C14.tex
% ============================================================================
% Chapitre de synthèse : aucun fait nouveau majeur ; reprise des faits de la
% bible (P00 à C13) pour la BANQUE SANAGA, LiZbiZ et l'ANPP. Faits du
% corpus pour l'ANPP : 40 paiements contrôlés par la Chambre des comptes,
% 29 OUI, 6 NON, 5 NEK (4 agents payés après leur départ ; pièces « aux
% archives »). Faits nouveaux mineurs (reportés dans la bible) : lettre de
% la COBAC de janvier 2029 demandant l'état des mesures ; demande de
% PHARMACAM d'un dossier de renouvellement pour janvier 2030 ; audit de
% l'ANTIC à l'ANPP programmé en septembre 2028 (Chapitre 10).
% ============================================================================

\chapter{Études de Cas Intégratrices}
\label{chap:cas-integrateurs}
\chaptermark{Cas intégrateurs}

\epigraph{%
    « Une méthode se juge à la fin d'une histoire,\\
    pas au début d'un chapitre. »%
}{Principe du cours}

\niveauChapitre{Expert}{%
    Maîtrise du Chapitre~\ref{chap:audit-technique}, du Chapitre~\ref{chap:grc-pratique}%
}

\begin{boxobjectifs}
\begin{itemize}
    \item Retracer, pour chacun des trois cas, le chemin de l'incident ou du
          constat initial jusqu'au rapport final.
    \item Relier gouvernance, appréciation et quantification des risques,
          audit, référentiels et droit dans une même mission.
    \item Lire l'évolution d'un dispositif dans le temps par les trois
          mesures $\mu$, $U$ et $H$.
    \item Reconnaître et éviter les erreurs typiques de l'auditeur et du
          gestionnaire des risques.
    \item Préparer l'épreuve de synthèse du cours.
\end{itemize}
\end{boxobjectifs}

\bigskip

% ============================================================================
\section*{Avant de commencer : trois histoires, une méthode}
% ============================================================================

Les treize chapitres précédents ont suivi trois organisations, chapitre
après chapitre\footnote{Comme dans tout le cours, les organisations, les
personnes, les dates et les montants sont fictifs~; les méthodes, les
normes et les textes cités sont exacts.}. La BANQUE SANAGA a été compromise
malgré son certificat, puis a reconstruit sa sécurité et sa gouvernance.
LiZbiZ a découvert des fraudes, puis a dû convaincre un client exigeant.
L'ANPP, organisme public, a dû répondre à la fois à un bailleur, à la
Chambre des comptes et à l'ANTIC.

\medskip

Ce chapitre ne présente aucune notion nouvelle. Il rassemble chaque
histoire en un seul fil, montre ce que chaque méthode y a apporté, en tire
les erreurs à ne pas commettre, puis propose trois études de cas complètes.
C'est le chapitre à relire avant l'épreuve de synthèse.

\separateur

% ============================================================================
\section{Une Méthode en Dix Questions}
\label{sec:c14-methode}
% ============================================================================

Toutes les missions du cours, quel que soit leur objet, répondent aux mêmes
dix questions, dans cet ordre.

\begin{tableaucours}{cL{6.2cm}L{3.4cm}c}{%
    \tableauHeader{N°} &
    \tableauHeader{Question} &
    \tableauHeader{Outil} &
    \tableauHeader{Chap.}
}
1 & Qui décide, qui répond, qui vérifie~? & RACI, trois lignes & \ref{chap:grc} \\
2 & Quel est le périmètre, et qu'en exclut-on~? & Périmètre défendable &
\ref{chap:smsi} \\
3 & Quels scénarios, avec quelle base de preuve~? & Fiche de risque, EBIOS &
\ref{chap:risque}, \ref{chap:ebios} \\
4 & Combien cela coûterait-il~? & FAIR, Monte-Carlo & \ref{chap:fair} \\
5 & Que décide-t-on, et qui accepte le reste~? & Plan, critères
d'acceptation & \ref{chap:iso27005} \\
6 & Quels textes s'imposent, lesquels sont choisis~? & Registre des
obligations & \ref{chap:referentiels}, \ref{chap:droit} \\
7 & Les mesures sont-elles en place et efficaces~? & Audit, exports figés &
\ref{chap:audit}, \ref{chap:audit-technique} \\
8 & Que peut-on conclure, et que ne sait-on pas~? & Décomposition, $\mu$,
$U$, $H$ & \ref{chap:conclure} \\
9 & Le système se corrige-t-il seul~? & Audit interne, revue de direction &
\ref{chap:smsi} \\
10 & Une seule vue pour la gouvernance~? & Matrice, tableau de bord &
\ref{chap:grc-pratique} \\
\end{tableaucours}
\vspace{-0.8em}
\begin{center}{\small\itshape Tableau~14.1. La méthode du cours en dix
questions}\end{center}

% ============================================================================
\section{BANQUE SANAGA : du Rançongiciel au Certificat}
\label{sec:c14-sanaga}
% ============================================================================

\subsection{Le fil}

\begin{tableaucours}{L{2.2cm}L{6.6cm}c}{%
    \tableauHeader{Date} &
    \tableauHeader{Événement} &
    \tableauHeader{Chap.}
}
Juin 2025 & Certificat sur le seul centre de données & \ref{chap:prologue} \\
Sept. 2026 & Rançongiciel par l'agence de Ngaoundéré~: 180 postes, 9
agences, 410~M FCFA & \ref{chap:prologue} \\
Oct. 2026 & « Qui a décidé du périmètre~? »~: gouvernance 2/4/2 &
\ref{chap:grc} \\
Nov. 2026 & R-AG-01 coté 16 puis 12~; registre 2025 revu~: 6/2/4 &
\ref{chap:risque} \\
Déc. 2026 & Critères d'acceptation, plan A1--A8 (193~M) & \ref{chap:iso27005} \\
Mars 2027 & Quantification~: 165~M par an, puis 19 après traitement &
\ref{chap:fair} \\
Sept. 2027 & Suivi du plan~: 5/2/3 & \ref{chap:iso27005} \\
Mars 2028 & Audit des agences~: 3/3/2~; Kribi, GAB & \ref{chap:audit},
\ref{chap:conclure} \\
Avr. 2028 & Déclaration d'applicabilité confrontée aux preuves~: 3/3/3 &
\ref{chap:iso27000} \\
Juin 2028 & Étape~2~: une majeure (GAB non appréciés) & \ref{chap:smsi} \\
Sept. 2028 & Certificat renouvelé sur toute la banque & \ref{chap:smsi} \\
Avr. 2029 & Page GRC du comité~: obligations 7/2/3 & \ref{chap:grc-pratique} \\
\end{tableaucours}
\vspace{-0.8em}
\begin{center}{\small\itshape Tableau~14.2. La BANQUE SANAGA, de 2025 à
2029}\end{center}

\subsection{Ce que les mesures racontent}

\begin{boxlecture}{Le risque R-AG-01 et la connaissance qu'on en a}
\textbf{Constatation.} La criticité résiduelle de R-AG-01 passe de 12
(2025) à 9 (mars 2027), puis à 6 (septembre 2027), confirmé par l'audit de
mars 2028. Sur la même période, les évaluations successives donnent~:
\begin{itemize}
    \item agences, prologue~: $\mu = 0{,}20$, $U = 0{,}40$~;
    \item agences, audit de mars 2028~: $\mu = 0{,}375$, $U = 0{,}25$~;
    \item obligations, avril 2029~: $\mu \approx 0{,}58$, $U = 0{,}25$.
\end{itemize}

\textbf{Analyse} (confiance moyenne~: les ensembles mesurés ne sont pas
identiques). Le risque technique a baissé plus vite que l'incertitude. Les
barrières techniques ont été posées et prouvées en dix-huit mois~; la part
d'inconnu, elle, est restée autour d'un quart, parce qu'elle s'est déplacée~:
des agences vers le personnel, puis vers les prestataires et les textes. La
banque n'a pas cessé d'ignorer~; elle a appris \emph{où} elle ignorait.
\end{boxlecture}

\subsection{Ce que le cas enseigne}

\begin{itemize}
    \item Un certificat ne protège que son périmètre~; le périmètre est une
          décision du conseil.
    \item Les défaillances de 2025 et de 2028 ont la même forme~: un risque
          connu, ou connaissable, laissé sans appréciation sur une partie du
          périmètre.
    \item Les mesures techniques se corrigent vite~; les mesures humaines et
          les relations avec les prestataires, lentement.
\end{itemize}

% ============================================================================
\section{LiZbiZ : de la Fraude au Client Exigeant}
\label{sec:c14-lizbiz}
% ============================================================================

\subsection{Le fil}

\begin{tableaucours}{L{2.2cm}L{6.6cm}c}{%
    \tableauHeader{Date} &
    \tableauHeader{Événement} &
    \tableauHeader{Chap.}
}
2026 & Écart de stock de 186~M FCFA~; fraude TRANSLOG (112,7~M)~; registre
R1--R8 (cours ACSI) & \ref{chap:risque} \\
2026 & Audit des quatre centres~: 8/9/3 & \ref{chap:audit-technique} \\
2026--2027 & Base de preuve du registre~: 4/1/3~; R8 recoté à 12 &
\ref{chap:risque} \\
Janv. 2027 & PHARMACAM exige une analyse~; étude EBIOS animée par
Mme~NKENG~; SO1 à 12 & \ref{chap:ebios} \\
2027 & Quantification simple de SO1~: 100~M par an & \ref{chap:fair} \\
2028 & Audit contractuel de PHARMACAM~; déclaration d'applicabilité 3/2/1 &
\ref{chap:audit}, \ref{chap:iso27000} \\
Oct. 2028 & Second passage technique~: 14/4/2~; règle de maintenance de
Garoua & \ref{chap:audit-technique} \\
2029 & Engagement d'un SMSI, certification visée fin 2029 & \ref{chap:smsi} \\
\end{tableaucours}
\vspace{-0.8em}
\begin{center}{\small\itshape Tableau~14.3. LiZbiZ, de 2026 à 2029}\end{center}

\subsection{Ce que le cas enseigne}

\begin{itemize}
    \item Une petite organisation n'a ni RSSI ni audit interne~: elle
          compense par la séparation des rôles qu'elle peut tenir
          (contrôleuse interne, auditeur externe indépendant) et par une
          hygiène de base vérifiable.
    \item Les fraudes internes et les attaques externes exploitent les
          mêmes faiblesses~; EBIOS les met sur le même plan.
    \item Le client est un régulateur de fait~: ses clauses (analyse,
          droit d'audit) font avancer la sécurité plus vite que la loi.
    \item Le défaut qui ouvre le scénario principal peut survivre à tous les
          plans s'il est «~temporaire~»~: la règle de Garoua l'a montré.
\end{itemize}

% ============================================================================
\section{ANPP : un Organisme Public sous Trois Regards}
\label{sec:c14-anpp}
% ============================================================================

\subsection{Le fil}

\begin{tableaucours}{L{2.2cm}L{6.6cm}c}{%
    \tableauHeader{Date} &
    \tableauHeader{Événement} &
    \tableauHeader{Chap.}
}
2026 & Chambre des comptes~: 40 paiements contrôlés, 29~\OUI{}, 6~\NON{},
5~\NEK{}~; quatre agents payés après leur départ (cours ACSI) &
\ref{chap:referentiels} \\
Mars 2027 & Un bailleur exige un profil NIST CSF~; base commune B1--B10~:
4/3/3 & \ref{chap:referentiels} \\
2028 & L'ANTIC annonce un audit de sécurité~; obligations légales~: 0/3/2 &
\ref{chap:droit} \\
2028--2029 & Tableau de bord GRC pour le conseil d'administration &
\ref{chap:grc-pratique} \\
\end{tableaucours}
\vspace{-0.8em}
\begin{center}{\small\itshape Tableau~14.4. L'ANPP, de 2026 à 2029}\end{center}

\begin{boxlecture}{Les paiements de l'ANPP}
\textbf{Constatation.} 29~\OUI{}, 6~\NON{}, 5~\NEK{} sur 40 paiements~:
$\mu \approx 0{,}73$~; $U = 0{,}125$~;
\[
H = -(0{,}725\log_2 0{,}725 + 0{,}15\log_2 0{,}15 + 0{,}125\log_2 0{,}125)
  \approx 0{,}336 + 0{,}411 + 0{,}375 \approx 1{,}12~\text{bit}.
\]

\textbf{Analyse} (confiance moyenne). Lu seul, le contrôle des paiements
paraît plutôt conforme. Rapproché de la base commune (B2~: comptes des
partants actifs, \NON{}) et des obligations légales (0/3/2), il change de
sens~: les quatre agents payés après leur départ sont le symptôme
financier d'un défaut de sécurité. Le même défaut apparaît dans trois
rapports, à trois destinataires~; il n'appelle qu'une correction.

\medskip
\hfill\lecture{0{,}725}{0{,}125}{1{,}12}
\end{boxlecture}

\subsection{Ce que le cas enseigne}

\begin{itemize}
    \item Un organisme public répond à plusieurs autorités, avec des textes
          parfois ambigus~; l'auditeur signale l'ambiguïté au lieu de la
          trancher.
    \item Une base commune de contrôles, testée une fois, évite trois
          chantiers et révèle les défauts communs.
    \item La sécurité et la régularité financière se rejoignent dans les
          contrôles d'accès.
\end{itemize}

% ============================================================================
\section{Dix Erreurs à ne pas Commettre}
\label{sec:c14-erreurs}
% ============================================================================

\begin{tableaucours}{cL{6.0cm}L{3.8cm}c}{%
    \tableauHeader{N°} &
    \tableauHeader{Erreur} &
    \tableauHeader{Où le cours l'a montrée} &
    \tableauHeader{Chap.}
}
1 & Prendre un certificat pour une preuve de sécurité & Hall de la BANQUE
SANAGA & \ref{chap:prologue} \\
2 & Lire un risque absent du registre comme un risque faible & Registre
2025 & \ref{chap:prologue} \\
3 & Laisser une décision de risque sous l'instance qui en répond &
Périmètre signé par le seul DG & \ref{chap:grc} \\
4 & Coter une vraisemblance sur un silence non surveillé & « Jamais
attaqués » & \ref{chap:risque} \\
5 & Réduire un risque par une mesure non testée & Cloisonnement de mars
2027 & \ref{chap:risque} \\
6 & Présenter un chiffre unique au lieu d'une distribution & 143,5~M &
\ref{chap:fair} \\
7 & Coter \OUI{} sur une déclaration & Rapport « tout est en place » &
\ref{chap:audit} \\
8 & Taire une limitation de l'étendue & GAB, prestataire & \ref{chap:audit} \\
9 & Résumer un audit par un pourcentage & « 37~\% » & \ref{chap:conclure} \\
10 & Reporter une conformité par un lien partiel & « Conforme ISO, donc
NIST » & \ref{chap:referentiels} \\
\end{tableaucours}
\vspace{-0.8em}
\begin{center}{\small\itshape Tableau~14.5. Les dix erreurs typiques}\end{center}

Toutes ont la même racine~: confondre ce que l'on sait et ce que l'on
suppose. Le cours n'a enseigné qu'une seule discipline, sous plusieurs
noms~: dire, à chaque étape, ce qui est établi, ce qui est contredit, et ce
que l'on ignore.

% ============================================================================
\section{Préparer l'Épreuve de Synthèse}
\label{sec:c14-epreuve}
% ============================================================================

L'épreuve de synthèse prend la forme d'une étude de cas, sur le modèle des
trois études ci-dessous~: un dossier (faits, extraits de documents,
résultats de tests) et des questions qui suivent la méthode en dix
questions. Le barème détaillé, la durée et les documents autorisés sont
fixés par l'enseignant et communiqués avec le calendrier de l'UE.

\begin{boxprocedure}
\textbf{Répondre à une étude de cas.}
\begin{itemize}
    \item Lire tout le dossier avant de répondre~; repérer les acteurs, le
          périmètre et ce qui est demandé.
    \item Séparer, pour chaque réponse, la constatation (le fait, la preuve)
          de l'analyse (l'interprétation, avec un niveau de confiance).
    \item Poser les calculs~: cotes, criticités, $\mu$, $U$ et $H$, avec les
          valeurs intermédiaires.
    \item Nommer ce que l'on ne sait pas, et dire comment le savoir.
    \item Citer les normes et les textes avec leur degré de vérification~;
          ne jamais inventer un numéro d'article.
    \item Conclure par une décision, et dire qui la prend.
\end{itemize}
\end{boxprocedure}

% ============================================================================
\section{Le Parcours de chaque Métier}
\label{sec:c14-metiers}
% ============================================================================

\begin{boxmetier}{RSSI}
\textbf{De Mme ABENA en 2026 à Mme ABENA en 2029}

Elle a obtenu le droit d'être entendue par le comité, cessé d'administrer ce
qu'elle surveille, remplacé ses déclarations par des preuves, et présenté
un indicateur défavorable (11~\% de clics) sans l'enjoliver.
\end{boxmetier}

\begin{boxmetier}{Auditeur en cybersécurité}
\textbf{Mme EKAMBI, M. BEYALA, CERTIQUAL AFRIQUE}

Trois auditeurs, trois postures, une même règle~: la preuve que l'on
obtient soi-même, la limitation que l'on déclare, le constat que l'on
maintient sous la pression, l'indépendance que l'on protège.
\end{boxmetier}

\begin{boxmetier}{Gestionnaire des risques cyber}
\textbf{Du registre au scénario}

Il est passé d'une liste cotée sans preuve à des scénarios décomposés,
reliés aux attaquants et à l'écosystème, chiffrés en francs CFA, avec une
base de preuve pour chaque cote.
\end{boxmetier}

\begin{boxmetier}{Responsable GRC}
\textbf{M. ATANGANA}

Il a relié quatre rapports en une page, et fait apparaître ce que le comité
ignorait. Son outil principal n'est pas un logiciel~: ce sont les liens
entre risques, obligations, contrôles et preuves.
\end{boxmetier}

\begin{boxmetier}{Consultant en mise en œuvre de SMSI}
\textbf{Le consultant de 2025, et celui de LiZbiZ}

Le premier a livré un certificat sur un périmètre commode~; celui que
LiZbiZ retiendra devra livrer un processus qui fonctionne sans lui, et un
périmètre qui inclut ses interfaces.
\end{boxmetier}

\separateur

\begin{boxresume}
\textbf{Ce qu'il faut retenir de ce chapitre~:}
\begin{itemize}
    \item \textbf{Dix questions} structurent toute mission~: rôles,
          périmètre, scénarios, chiffrage, décisions, textes, preuves,
          conclusion, amélioration, vue d'ensemble.
    \item \textbf{BANQUE SANAGA}~: le risque baisse vite, l'inconnu se
          déplace~; les échecs de 2025 et de 2028 ont la même forme.
    \item \textbf{LiZbiZ}~: fraude et attaque exploitent les mêmes
          faiblesses~; le client fait avancer la sécurité.
    \item \textbf{ANPP}~: un même défaut d'accès apparaît dans trois
          rapports~; une base commune le corrige une fois.
    \item \textbf{Dix erreurs}, une racine~: confondre ce que l'on sait et
          ce que l'on suppose.
\end{itemize}
\end{boxresume}

% ============================================================================
\section*{Exercices}
% ============================================================================

\begin{boxexo}[title={\ding{45}\quad Étude de cas~14.1. La réponse à la COBAC \niveauexpert}]
En janvier 2029, la COBAC demande à la BANQUE SANAGA «~l'état des mesures
prises depuis l'incident de septembre 2026~».
\begin{enumerate}[(a)]
    \item Identifiez le cadre de la demande et le degré de vérification du
          texte (Chapitre~\ref{chap:droit}).
    \item Rédigez le plan de la réponse, en suivant la méthode en dix
          questions.
    \item Rédigez la partie «~risque~» de la réponse~: évolution de la
          criticité de R-AG-01, avec la base de preuve de chaque étape, et
          la quantification avant et après traitement.
    \item Rédigez la partie «~ce que la banque ne sait pas encore~», et le
          calendrier pour le savoir.
    \item Qui signe la réponse, et pourquoi~?
\end{enumerate}
\end{boxexo}

\begin{boxexo}[title={\ding{45}\quad Étude de cas~14.2. Le dossier de renouvellement de PHARMACAM \niveauexpert}]
PHARMACAM demande à LiZbiZ, pour le renouvellement de janvier 2030, un
dossier de sécurité.
\begin{enumerate}[(a)]
    \item Listez les pièces du dossier, en distinguant ce qui est prouvé et
          ce qui est déclaré.
    \item Retracez l'évolution du scénario SO1 de 2027 à 2028 (criticité,
          phases couvertes, preuves).
    \item Comparez les deux passages de l'audit technique (8/9/3 et 14/4/2)
          centre par centre, et rédigez le paragraphe qui les présente à
          PHARMACAM sans moyenne trompeuse.
    \item Quelle décision recommandez-vous à M.~ETEME sur la certification
          ISO/IEC~27001, et à quelle échéance~?
\end{enumerate}
\end{boxexo}

\begin{boxexo}[title={\ding{45}\quad Étude de cas~14.3. Un défaut, trois rapports \niveauexpert}]
À l'ANPP, le défaut de désactivation des comptes des agents partis apparaît
dans trois travaux~: le contrôle des paiements de la Chambre des comptes,
la base commune (B2), et les obligations légales (article~29, alinéa~2).
\begin{enumerate}[(a)]
    \item Montrez, pour chacun des trois destinataires (M.~NJOYA, le
          bailleur, l'ANTIC), ce que le défaut signifie dans sa langue.
    \item Rédigez une seule recommandation, avec sa preuve de mise en œuvre,
          qui réponde aux trois.
    \item Proposez le test automatique qui permettra de vérifier chaque
          nuit que le défaut ne revient pas, et dites comment l'audit le
          validera.
    \item Recalculez $\mu$, $U$ et $H$ de la base commune si B2 devient \OUI{}
          et si B4 et B6 sont établis (\OUI{} et \NON{}).
\end{enumerate}
\end{boxexo}

\bigskip

\begin{boxinfo}
\textbf{Transition.} Les trois cas reposent sur des systèmes que leurs
propriétaires détiennent. Les risques qui viennent sont ailleurs~: chez les
prestataires et dans le nuage, dans l'intelligence artificielle, dans
l'audit qui devient continu. La Partie~V les aborde au
Chapitre~\ref{chap:emergents}.
\end{boxinfo}
